% Expository adaptation, October 7, 2026. Original mathematics: OpenAI, September 23, 2026.
% Reorganized and expanded for pedagogy; see README.md for the pinned source and scope.
\section{Computing numerical values one digit at a time}
\label{sec:arith}

We have defined the estimator as a mathematical function of a finite random environment. Its values involve arithmetic, conditional averages, clipping, and order statistics. This section explains how finite numerical queries evaluate those values without keeping a long approximation in every caller.

There are two notions of accuracy. The statistical budget determines which estimator value is being computed. The numerical digit place determines how accurately that fixed value is read. Increasing the latter does not improve the former, and it must not change the list or port choices that define the estimator. We therefore give each fixed task value an exact signed-digit representation and implement a request for any one digit by a terminating finite calculation.

Why return one digit at a time? If a multiplication first computed and saved a whole approximation to its left operand, that word would remain live while the right operand made recursive calls. A digit interface lets the arithmetic circuit request only the bounded pieces it needs and recompute them when necessary.

A query returns one bounded digit.  During its evaluation, the space
needed by a suspended caller must be bounded by the size of its
structural descriptors and a constant times the decrease from its
output clock to the requested operand clock.  An operand clock may
exceed the output clock by an amount proportional to descriptor size;
the structural term also pays for this overshoot.  These differences
will telescope when arithmetic queries are nested inside graph procedures.
Thus the relevant question is which data remain live at each operand request.

\subsection{Exact values from bounded signed digits}

We use integer clocks $t,j\geq 0$.  Known binary scales are separated
from numerical values, so a normalized value belongs to $[-1,1]$.
The fixed digit alphabet is
\[
 \mathcal A=\{-6,-5,\ldots,5,6\}.
\]
Every digit is encoded by a fixed number of bits.  An operand stream
$(x_j)_{j\geq0}$ denotes the exact real number
\begin{equation}
 x=\sum_{j\geq0}x_j2^{-j},\qquad x_j\in\mathcal A.
 \label{eq:arith-stream}
\end{equation}
In particular, for its truncation $x^{[K]}=\sum_{j=0}^Kx_j2^{-j}$,
\begin{equation}
 |x-x^{[K]}|\leq 6\,2^{-K}.
 \label{eq:arith-tail}
\end{equation}
The stream need not be a usual binary expansion, and no uniqueness
of the expansion is required.
Redundant signed-digit representations have a long history in
arithmetic~\cite{Avizienis1963}.  The redundancy lets consecutive
finite approximations define an exact stream even when they use
different truncations.

An arithmetic root has public parameters, its clock, and a description
of its operand interfaces.  We require these parameters to be
regenerable from the outer controller records without operand or graph
queries.  They are fixed inputs to each numerical computation; their
full binary encodings are not copied at every gate.  An operand label
specifies its interface and descriptor; it never specifies an unknown
vertex address from a previous cursor position.

A finite approximation at place $t$ need not extend the approximation obtained at place $t-1$. Rather than enforce matching binary prefixes, take their scaled difference as the next signed digit. The resulting partial sums telescope to the exact value. The error bound keeps each difference in a fixed finite alphabet, so only a fixed number of residue bits is needed to recover it.

\begin{lemma}[Exact streams from approximation residues]
\label{lem:arith-residue-stream}
Suppose $|z|\leq1$ and, for each $t\geq0$, a deterministic integer
$N_t(z)$ satisfies
\begin{equation}
 |N_t(z)-2^tz|<2.
 \label{eq:arith-approximation}
\end{equation}
Set
\begin{equation}
 z_0=N_0(z),\qquad z_t=N_t(z)-2N_{t-1}(z)\quad(t\geq1).
 \label{eq:arith-difference-digit}
\end{equation}
Then $z_t\in\mathcal A$ and these digits represent $z$ exactly.
For any fixed integer $q\geq4$, each output digit is determined by
the required $N_t$ residues modulo $2^q$.
\end{lemma}

\begin{proof}
There are three things to check: the digits remain bounded, their weighted sum has the correct limit, and their residues determine them uniquely. Writing $e_t=N_t-2^tz$ gives
$z_t=e_t-2e_{t-1}$ for $t>0$, hence $|z_t|<6$; the weaker alphabet
bound $|z_t|\leq6$ is convenient throughout.  At place zero,
$|N_0|<3$, so $|N_0|\leq2$.  The partial sum through place $T$ is
$2^{-T}N_T$, which tends to $z$ by Equation~\eqref{eq:arith-approximation}.
Finally, reduction modulo $2^q$ is injective on the possible signed
digits because $2^q>12$.  Thus the signed difference can be recovered
from the two residues without forming the full integers $N_t$.
Concretely, for $t>0$ compute
$(N_t\bmod 2^q)-2(N_{t-1}\bmod 2^q)$ modulo $2^q$, then use the
fixed finite table which sends each occurring residue to its unique
representative in $\mathcal A$.  At $t=0$, decode the residue of
$N_0$ in the same way.  This final decoding has constant storage and
introduces no comparison of unknown exact real numbers.
\end{proof}

The integer $N_t$ is always defined before taking its residue.
Different places may use different truncations or finite circuits.
Only determinism and Equation~\eqref{eq:arith-approximation} are
required; consistency between the finite approximations is supplied
by Equation~\eqref{eq:arith-difference-digit}.  In particular, when a
digit at $t$ uses $N_{t-1}$, the latter is evaluated by its own
place-$(t-1)$ definition.

We compute the required residues by finite Boolean circuits.  To bound
what survives an operand query, we evaluate those circuits from paths
of child choices and regenerate their absolute addresses as needed.
Here $d\geq1$ measures the size of the immediate structural descriptors;
in the application it is the current group's slack.

The local record of a suspended numerical query contains every local bit
that must be retained across an operand query: its binary gate path,
return modes, pending bounded results, stored descriptor choices, and
any explicitly specified retained word.  The input and master choice,
the common vertex cursor, and one reusable scratch area are shared
facilities.  Temporary gate addresses, digit places and scan counters
may occupy scratch while they are being computed.  Any such data needed
after a recursive graph query must either be in the charged local
record or be reconstructible from the root parameters and that record.
An operand-dependent word retained across the query is charged even if
it was originally formed in scratch.

In particular, reconstructing an absolute digit place may use a binary
counter in scratch, while the suspended record stores only the choices
needed to recover that place.  The next lemma makes this distinction
precise for the finite circuits used below.

A second source of excess memory is saving an absolute circuit-gate address at every level. Instead, retain the finite choices along the route from the output gate. The absolute address of the current gate can be regenerated in shared scratch from that route. This trades recomputation for a smaller retained record.

\begin{lemma}[Evaluation from regenerated gate addresses]
\label{lem:arith-gate-evaluation}
Consider a finite Boolean circuit of bounded fan-in, with a specified
output gate.  Conditional gates may evaluate their selector first and
then only their selected branch.  Suppose the type and immediate predecessors of a gate,
or its operand label and digit place if it is an input gate, can be
computed from the root parameters and a path of child choices from
the output gate.  Suppose these computations use $O(B+t+d)$ scratch
bits and make no operand or graph queries.  Then depth-first
evaluation retains $O(1)$ bits per gate on its current path, and
requires no saved absolute gate address per gate.  Its suspended
record above an operand input has length at most a constant times
the length of that path, plus any explicitly retained root-local
record.
\end{lemma}

\begin{proof}
At a gate, save its next-child or return mode and the already computed
Boolean input, if any.  At a conditional gate also save the selected
branch bit.  The bounded fan-in makes this a constant-size
record.  To locate the next child, scan the saved path from its root
and regenerate successive gate descriptions in scratch.  The final
description determines the next child choice.  On return, its Boolean
result is combined with the saved bit and the local record is popped.
All temporary addresses and scan counters are discarded before an
operand query; a later scan recreates them from the same root and
saved choices.  Circuit acyclicity gives termination.  This is an
evaluation of the circuit unfolded into a formula; repeated gates are
recomputed, so no array of previously computed values is stored.
\end{proof}


The next theorem states exactly what the graph computation will need from arithmetic. The elementary operations, multiplication, reciprocals, and indexed reductions are proved in the appendix. Their proofs track what remains stored at a recursive operand request; a bound on temporary arithmetic workspace alone would not be enough.

\subsection{The numerical interface}


We state the composition result with its syntactic hypotheses.
They concern finite arithmetic recipes and do not assert any
property of the underlying graph operations.

Application to graph procedures requires a separate certificate.  The
semantic descriptor domains and structural validity conventions are
fixed independently of the requested numerical clock; selected operands
already have defined digit interfaces; reciprocal arguments lie in the
stipulated positive intervals; and root parameters remain recoverable
after lower calls.  Clock-dependent truncations and comparison circuits
may change $N_t$, but not its target recipe value.
Lemmas~\ref{lem:ctrl-recipes} and~\ref{lem:ctrl-replay} verify these
conditions for the estimator procedures.  An incoming rank over the
whole vertex universe has no prescribed $2^{O(d)}$ descriptor bound and
uses the separate construction of Lemma~\ref{lem:ctrl-score-rank}.

\begin{theorem}[Uniform numerical interface]
\label{thm:arith-interface}
Fix the digit alphabet, positive reciprocal intervals, and
constants bounding the following recipe class.  A recipe at
slack $d\geq1$ has a fixed number of stages built from bounded
arithmetic operations, fixed piecewise formulas using maxima,
minima and clipping, the reciprocals of
Lemma~\ref{lem:arith-reciprocal}, indexed ranks as in
Lemma~\ref{lem:arith-rank}, and sums and known scales as in
Lemma~\ref{lem:arith-sums-scales}.  The stage and descriptor
descriptions are generated from the root parameters by index
computations using $O(B+d+t)$ scratch and no operand queries, and the
total non-reduction dependency depth is bounded independently
of $d$.  Each indexed domain has size $2^{O(d)}$ and descriptors
of length $O(d)$.  Discrete domain and validity decisions are
structural leaves or buffered tests with their own fixed
precision conventions.  Public coefficient leaves have exact
normalized digit procedures using $O(B+d+t)$ shared scratch
and making no operand or recursive graph queries.

Then each normalized recipe value has a deterministic exact
$\mathcal A$-stream.  For its digit at clock $t$ there are uniform
constants $a,b,c$ such that every requested operand digit $j$
satisfies
\begin{equation}
 j\leq t+cd,\qquad
 \text{total suspended local binary record length}
       \leq ad+b(t-j).
 \label{eq:arith-export}
\end{equation}
An output code bit satisfies the same bounds.  A structural leaf
is requested at clock zero and has suspended local length at
most $ad+bt$.  The maximum intrinsic local record length is
$O(d+t)$.

All gate, digit and descriptor addresses can be regenerated from
root parameters and the saved local records using
$O(B+d+t)$ shared scratch and no recursive graph calls.  When
$d,t=O(B)$ this is $O(B)$ scratch.  If a public rational
coefficient has $O(B)$-bit numerator and denominator that can
be generated in $O(B)$ scratch, its queried digits in this clock
range also require only that shared $O(B)$ scratch.
\end{theorem}

\begin{proof}
The primitive constructions in the appendix define deterministic integers $N_t$.
Lemma~\ref{lem:arith-residue-stream}
turns them into exact streams.  When composing primitives,
their operands are these exact streams; numerical approximation
errors are not accumulated along the recipe as if they were
statistical errors.  Instead, each primitive's deterministic
$N_t$ is an approximation to its own exact mathematical value.
The construction uses only residues of $N_t$ and $N_{t-1}$ and
a constant number of output code bits, which adds fixed
overshoot and fixed local overhead.

For the record bound, choose one common positive clock
coefficient $b$ for the finite primitive and stage types.  In any
bound $L\leq a_0d+b_0(t-j)$ with $j-t\leq c_0d$, replacing
$b_0$ by $b\geq b_0$ is justified by replacing $a_0$ by
$a_0+(b-b_0)c_0$.  Thus the same $b$ can be used everywhere.
There are a bounded number of stages, each with overhead $O(d)$;
within a balanced reduction the elementary overheads have total
$O(d)$ by Lemma~\ref{lem:arith-sums-scales}.  Along any path
the intermediate clock differences telescope, giving
Equation~\eqref{eq:arith-export} and total overshoot $O(d)$.
The same sum, terminated at an internal arithmetic operation
rather than an operand call, gives the intrinsic bound explicitly
as follows.  If that operation has current clock $u$, its own
intrinsic record has length at most $C(d+u)$, while the outer
saved portion has length at most $A d+b(t-u)$.  Choose the
common $b$ also to satisfy $b\geq C$.  Their sum is then at most
$(A+C)d+bt$, because $u\geq0$.  Increasing the stage constants
as above makes this same choice of $b$ valid in all local bounds.
The early-word case is already included in
Lemma~\ref{lem:arith-reciprocal}: those retained bits are part of
the local record and are paid for at the actual original-input
clock by Equation~\eqref{eq:arith-early-word-payment}.

Structural validity branches are finite branches in this circuit
organization, with length at most its full intrinsic bound.
They ask the structural operand at clock zero.  Buffered tests
are separate recipes with structural starting clock zero and
precision fixed by Lemma~\ref{lem:arith-buffered-test}; they are
not altered when a surrounding value is requested more accurately.
Thus no discontinuous comparison of exact reals at unbounded
precision is part of the interface.

All circuit constructions in the appendix use padded domains and explicit index rules.  A saved gate choice selects an operand, a subtree,
a bit of a known-width word, or one of the finite primitive modes.
Scanning these choices from the root recovers the current
descriptor and digit places by integer arithmetic.  The
reciprocal's extra records additionally specify its early-prefix
formation state.  None of these computations refers to an
ancestor cursor address or requires an operand value to discover
the next gate.  Their integers have $O(B+d+t)$ bits; repeated
scans reuse one scratch area.  Lemma~\ref{lem:arith-gate-evaluation}
therefore applies at every Boolean portion of the recipe.

Finally, for a public coefficient $P/Q$, $Q>0$, producing a
place-$r$ approximation uses integer division of $2^rP$ by $Q$,
with signed floors when needed.  The inputs and intermediate
remainders have $O(B+r)$ bits.  Standard finite binary division
and multiplication use that much space, and extracting the
low residue or difference digit requires no larger word.
All operand bits in this computation are public, so it makes no
recursive graph call and can use shared scratch throughout.
For $r=O(B)$ its scratch bound is $O(B)$.  In contrast, any word
whose bits depend on recursive operands and must survive a new
operand request has already been charged to the local record.
\end{proof}

\begin{remark}[Use of clocks]
\label{rem:arith-clocks}
Theorem~\ref{thm:arith-interface} is a finite-request statement
for every $t$.  Its exact-stream conclusion follows by considering
these finite requests separately, and does not require a machine
to choose structural domains or port orders at infinite precision.
For the actual top-level computation all reached clocks will be
$O(B)$ by the allowance calculation in the next section.  Known
rate and field coefficients have $O(B)$-bit descriptions, and a
fixed rational generator probability raised to a walk length
$O(d)$ has numerator and denominator of $O(d)$ bits; these are
instances of the public-coefficient provision above.
\end{remark}

